\section{Introduction}

Physical networks---such as neuronal arbors, vascular systems, plant roots, and fungal hyphae---are constrained not only by their connectivity, but by the material and geometric conditions under which that connectivity is realized. For decades, these systems have been studied through optimization-based principles, most prominently wiring-length minimization and its refinements, which model networks as collections of one-dimensional links optimized for global efficiency.

Recent advances in three-dimensional imaging and reconstruction have made it possible to test these principles at unprecedented resolution. A key outcome of this empirical progress has been the recognition that classical wiring-economy models are insufficient to account for several robust and ubiquitous local features of real physical networks, including higher-degree junctions, non-planar branching, broad angle distributions, and orthogonal sprouting.

In response, surface-minimization models have been introduced that replace one-dimensional links with smooth, two-dimensional manifolds embedded in three-dimensional space. In these models, physical networks are represented as collections of surface ``sleeves'' joined smoothly at junctions, and network morphology is determined by minimizing an area functional subject to a non-collapse (thickness or capacity) constraint. Formulated in this way, the problem admits a natural interpretation in terms of classical minimal surfaces, closely related to the Nambu--Goto action for string worldsheets.\cite{Meng2026SurfaceOptimization}
The analogy is structural rather than ontological: the point is the shared mathematical form of a single classical extremum within a fixed topology class, not an attempt to import quantum string dynamics into biological modeling.

These surface-minimization approaches have achieved notable success. They correctly predict several local branching motifs that eluded earlier theories, including stable trifurcations, angle asymmetry regimes, and the prevalence of orthogonal sprouts across diverse biological and physical systems. Earlier transport-network optimization frameworks successfully captured related geometric constraints in reduced settings, but lacked the degrees of freedom required to reproduce these higher-order branching motifs.\cite{Katifori2010,Ronellenfitsch2015} Importantly, these predictions are obtained without system-specific tuning, suggesting that local geometric constraints play a fundamental role in shaping physical networks. Although Meng et al. acknowledge dynamical and developmental influences, their framework remains a classical saddle-point description and therefore inherits the structural limits identified here.

Despite these successes, the theoretical status of surface-minimization models remains unclear. While their empirical accuracy at the level of local motifs is well established, it is not evident which aspects of network organization they can determine in principle, and which lie beyond their scope. In particular, questions concerning topology selection, the admissibility of smooth continuation under competing constraints, and consistency across scales are often discussed informally but are not resolved within the surface-minimization framework itself.

The purpose of this paper is to clarify this boundary. We show that surface-minimization network models can be understood as classical saddle-point solutions of a constrained geometric action, evaluated within a restricted topology class appropriate to local branching. When viewed in this way, their explanatory power---and their limitations---become sharply defined. Certain classes of questions are resolved by the classical saddle, while others are structurally undecidable without introducing an additional principle.

Our central claim is that surface minimization, while necessary and highly informative, is generically incomplete as a theory of physical network organization. This incompleteness is not a matter of missing parameters or insufficient optimization, but a direct consequence of reliance on a single classical solution under fixed constraints. We argue that resolving the resulting ambiguities requires an explicit notion of admissibility: a rule that determines when smooth continuation is permitted and when topological reorganization becomes necessary.

The paper proceeds as follows. We begin by clarifying the interpretation of surface-minimization network models as classical saddle-point solutions of a constrained geometric action, a perspective that is used throughout to assess both their explanatory power and their limitations. Section~\ref{sec:incompleteness} delineates three distinct sources of structural incompleteness inherent to classical minimal-surface descriptions. In Section~\ref{sec:admissibility}, we articulate a minimal admissibility principle required to resolve the resulting indeterminacies while remaining consistent with the empirical successes of surface minimization. We conclude with a brief discussion of implications and open directions.

Throughout, we adopt a conservative posture. We do not propose a replacement theory, introduce new empirical claims, or challenge the validity of existing surface-minimization results. We aim to complete the logical arc they initiate by identifying the precise boundary between classical geometric optimization and the broader class of phenomena governed by admissibility.

\begin{center}
\fbox{\parbox{0.92\linewidth}{
\textbf{Formal setting.}
We treat surface-minimization network models as classical variational problems that minimize an area functional over smooth embedded manifolds subject to capacity constraints, evaluated within a fixed topology (continuation) class. All results below concern the structural decision power of this classical saddle-point formulation.
}}
\end{center}
